% GENERATED FILE. Edit canonical lessons, then run npm run generate:books.
% canonical-source-hash: fnv1a64:eed5580d5d8353ad
% canonical-lessons: JA-C96-shikaku, JA-C96-sainou, JA-C96-waza, JA-C96-yoyuu, JA-C96-tsugou

\chapter{Qualification, Talent, Technique, Room to spare, Convenience}
\label{ch:ja-qualification-talent-technique-room-to-spare-convenience}
\hlchaptermodality{96}

\begin{chapteropening}
\textbf{By the end of this chapter:} \emph{I can say a qualification, talent, a technique, room to spare and convenience.}

Complete the last lesson of chapter 96: i can say a qualification, talent, a technique, room to spare and convenience.
\end{chapteropening}

\section[\texorpdfstring{\ja{しかく} (shikaku)}{しかく (shikaku)}]{\ja{しかく} (shikaku) — a qualification}

\label{lesson:JA-C96-shikaku}



\begin{quote}
Before the new one: say the Japanese for belongings, then the Japanese for a rich person.
\end{quote}

\subsection*{You'll want to know: \ja{しかく}}
\textbf{\ja{しかく}} — \emph{shikaku} — \textquotedblleft{}a qualification\textquotedblright{}.

\begin{grammarlens}[title={What you've built}]
One more word for this chapter.
\end{grammarlens}

\subsection*{Your turn}
\begin{itemize}
\raggedright
  \item \emph{Say it:} \emph{shikaku}
  \item \emph{Say it:} \emph{shikaku}, once more
  \item \emph{Say it:} say \emph{shikaku}
  \item \emph{Recall:} say the Japanese for belongings, then the Japanese for a rich person, then say \emph{shikaku} again
\end{itemize}

\begin{tcolorbox}[breakable,colback=teal!4,colframe=teal!35!black,title={Before you move on}]
What does \ja{しかく} mean? (\textquotedblleft{}A qualification\textquotedblright{}.) Say it once more.
\end{tcolorbox}
\section[\texorpdfstring{\ja{さいのう} (sainou)}{さいのう (sainou)}]{\ja{さいのう} (sainou) — talent}

\label{lesson:JA-C96-sainou}



\begin{quote}
Before the new one: say the Japanese for a rich person, then the Japanese for a qualification.
\end{quote}

\subsection*{You'll want to know: \ja{さいのう}}
\textbf{\ja{さいのう}} — \emph{sainou} — \textquotedblleft{}talent\textquotedblright{}.

\begin{grammarlens}[title={What you've built}]
One more word for this chapter.
\end{grammarlens}

\subsection*{Your turn}
\begin{itemize}
\raggedright
  \item \emph{Say it:} \emph{sainou}
  \item \emph{Say it:} \emph{sainou}, once more
  \item \emph{Say it:} say \emph{sainou}
  \item \emph{Recall:} say the Japanese for a rich person, then the Japanese for a qualification, then say \emph{sainou} again
\end{itemize}

\begin{tcolorbox}[breakable,colback=teal!4,colframe=teal!35!black,title={Before you move on}]
What does \ja{さいのう} mean? (\textquotedblleft{}Talent\textquotedblright{}.) Say it once more.
\end{tcolorbox}
\section[\texorpdfstring{\ja{わざ} (waza)}{わざ (waza)}]{\ja{わざ} (waza) — a technique}

\label{lesson:JA-C96-waza}



\begin{quote}
Before the new one: say the Japanese for a qualification, then the Japanese for talent.
\end{quote}

\subsection*{You'll want to know: \ja{わざ}}
\textbf{\ja{わざ}} — \emph{waza} — \textquotedblleft{}a technique\textquotedblright{}.

\begin{grammarlens}[title={What you've built}]
One more word for this chapter.
\end{grammarlens}

\subsection*{Your turn}
\begin{itemize}
\raggedright
  \item \emph{Say it:} \emph{waza}
  \item \emph{Say it:} \emph{waza}, once more
  \item \emph{Say it:} say \emph{waza}
  \item \emph{Recall:} say the Japanese for a qualification, then the Japanese for talent, then say \emph{waza} again
\end{itemize}

\begin{tcolorbox}[breakable,colback=teal!4,colframe=teal!35!black,title={Before you move on}]
What does \ja{わざ} mean? (\textquotedblleft{}A technique\textquotedblright{}.) Say it once more.
\end{tcolorbox}
\section[\texorpdfstring{\ja{よゆう} (yoyuu)}{よゆう (yoyuu)}]{\ja{よゆう} (yoyuu) — room to spare}

\label{lesson:JA-C96-yoyuu}



\begin{quote}
Before the new one: say the Japanese for talent, then the Japanese for a technique.
\end{quote}

\subsection*{You'll want to know: \ja{よゆう}}
\textbf{\ja{よゆう}} — \emph{yoyuu} — \textquotedblleft{}room to spare\textquotedblright{}.

\begin{grammarlens}[title={What you've built}]
One more word for this chapter.
\end{grammarlens}

\subsection*{Your turn}
\begin{itemize}
\raggedright
  \item \emph{Say it:} \emph{yoyuu}
  \item \emph{Say it:} \emph{yoyuu}, once more
  \item \emph{Say it:} say \emph{yoyuu}
  \item \emph{Recall:} say the Japanese for talent, then the Japanese for a technique, then say \emph{yoyuu} again
\end{itemize}

\begin{tcolorbox}[breakable,colback=teal!4,colframe=teal!35!black,title={Before you move on}]
What does \ja{よゆう} mean? (\textquotedblleft{}Room to spare\textquotedblright{}.) Say it once more.
\end{tcolorbox}
\section[\texorpdfstring{\ja{つごう} (tsugou)}{つごう (tsugou)}]{\ja{つごう} (tsugou) — convenience, availability}

\label{lesson:JA-C96-tsugou}



\begin{quote}
Before the new one: say the Japanese for a technique, then the Japanese for room to spare.
\end{quote}

\subsection*{You'll want to know: \ja{つごう}}
\textbf{\ja{つごう}} — \emph{tsugou} — \textquotedblleft{}convenience, availability\textquotedblright{}.

\begin{grammarlens}[title={What you've built}]
One more word for this chapter.
\end{grammarlens}

\subsection*{Your turn}
\begin{itemize}
\raggedright
  \item \emph{Say it:} \emph{tsugou}
  \item \emph{Say it:} \emph{tsugou}, once more
  \item \emph{Say it:} say \emph{tsugou}
  \item \emph{Recall:} say the Japanese for a technique, then the Japanese for room to spare, then say \emph{tsugou} again
\end{itemize}

\begin{tcolorbox}[breakable,colback=teal!4,colframe=teal!35!black,title={Before you move on}]
What does \ja{つごう} mean? (\textquotedblleft{}Convenience, availability\textquotedblright{}.) Say it once more.
\end{tcolorbox}
